% Lecture 04: Функции, ссылки и const

\section{Лекция 4. Функции, ссылки и const}
\label{sec:lecture-04}

В первых программах вся последовательность действий помещалась в
\texttt{main}. В утилите инвентаризации уже появились разбор аргументов,
проверка файлов, подсчёты и печать таблицы. Они работают, но спрятаны в одном
длинном блоке. Следующее изменение заставляет заново искать границы каждой
задачи и повышает риск случайно сломать соседнюю.

\begin{importantbox}[Главный вопрос]
Как превратить последовательность команд в набор именованных операций с
ясными контрактами: что операция получает, что изменяет и что возвращает?
\end{importantbox}

\subsection{Функция как единица проектирования}

Функция --- именованная операция. Она позволяет заменить повторённый алгоритм
одним вызовом, дать намерению имя и отделить интерфейс от реализации.

\begin{cppcode}
int square(int value) {
    return value * value;
}

const int area = square(6);
\end{cppcode}

Читающий \texttt{square(6)} сначала видит задачу, а детали умножения может
изучить отдельно. Исправление алгоритма делается в одном месте. Так функции
поддерживают читаемость, повторное использование, тестирование и декомпозицию.

Но механически выносить каждую строку не цель. Хорошая функция выполняет одну
понятную задачу, имеет ясные входы и выходы, разумное имя и уровень абстракции.
\texttt{doStuff()} скрывает смысл; цепочка функций по одной очевидной строке
заставляет читателя прыгать по файлу. Граница полезна, когда часть программы
можно назвать и понять как самостоятельную операцию.

\begin{definition}
\textbf{Интерфейс функции} на этом уровне --- имя, тип возвращаемого значения
и типы параметров. \textbf{Реализация} --- тело, которое выполняет обещанную
операцию.
\end{definition}

\subsection{Определение, параметры и вызов}

\begin{cppcode}
int power(int base, int exponent) {
    int result{1};
    for (int i = 0; i < exponent; ++i) {
        result *= base;
    }
    return result;
}

const int answer = power(3, 4);
\end{cppcode}

Разберём запись:
\textbf{\texttt{int}:} return type: вызов вычисляет значение этого типа; \textbf{\texttt{power}:} function name; \textbf{\texttt{(int base, int exponent)}:} parameter list; \textbf{фигурные скобки:} body --- выполняемые действия; \textbf{\texttt{return result;}:} завершает вызов и передаёт значение caller; \textbf{\texttt{power(3, 4)}:} call; \texttt{3} и \texttt{4} --- arguments.

Параметр --- локальное имя внутри definition, argument --- выражение в месте
call. Порядок важен: \texttt{power(3, 4)} и \texttt{power(4, 3)} связывают
разные значения с \texttt{base} и \texttt{exponent}. Компилятор проверяет
число аргументов, допустимые conversions и тип результата.

Runtime управление переходит в тело функции, затем возвращается в точку после
call. Каждый вызов получает собственные параметры и локальные объекты. Это
семантика языка. Отдельный stack frame и call stack --- полезная обычная
модель реализации, но стандарт не требует конкретного CPU stack или ABI.

Функция без значения результата имеет return type \texttt{void}:

\begin{cppcode}
void print_score(int score) {
    std::cout << "score=" << score << '\n';
}
\end{cppcode}

В \texttt{void}-функции допустим \texttt{return;} без expression. В функции с
результатом значение должно быть возвращено на каждом достигнутом пути
(кроме специальных правил для \texttt{main}). Код после безусловного
\texttt{return} недостижим.

\subsection{Declaration, prototype, definition и call}

Чтобы проверить call, compiler должен раньше увидеть declaration:

\begin{cppcode}
int power(int base, int exponent); // declaration, или prototype

int main() {
    return power(2, 5);             // call
}

int power(int base, int exponent) { // definition
    int result{1};
    for (int i = 0; i < exponent; ++i) {
        result *= base;
    }
    return result;
}
\end{cppcode}

Declaration сообщает имя и function type; завершающий semicolon означает,
что тела здесь нет. Definition --- тоже declaration, но дополнительно
задаёт body. Call не объявляет функцию. Имена параметров в prototype можно
опустить, но содержательные имена документируют роли.

В блоке~1 translation units компилировались независимо. Обычно общие
declarations помещают в header, а definition --- в \texttt{.cpp}, чтобы
каждая единица трансляции могла проверить вызов. Полные правила headers, ODR и
linkage будут развиваться позже; фундамент таков: compiler проверяет видимый
declaration, linker затем должен найти подходящее definition.

В строгой терминологии signature не всегда включает всё, что новичок называет
«заголовком функции». Здесь безопаснее говорить явно: function name,
parameter types и return type; перегрузка различается списком параметров, но
не одним return type.

\subsection{Return value как часть контракта}

Функция-вычисление естественно возвращает результат:

\begin{cppcode}
int absolute(int value) {
    if (value >= 0) {
        return value; // ранний выход упрощает основной путь
    }
    return -value;
}
\end{cppcode}

Ранний return полезен для простого guard condition; слишком много выходов
может, наоборот, затруднить обзор. Контракт должен объяснять смысл результата,
включая special values. \texttt{return} возвращает вычисленное значение,
поэтому обычный return-by-value безопасен после уничтожения local.

\begin{ubbox}[Dangling return]
Нельзя вернуть pointer или reference на обычный local object и затем получить
через него доступ: при выходе lifetime local заканчивается. Адресное значение
может сохраниться, но живого объекта по нему уже нет.
\end{ubbox}

\begin{cppcode}
int* bad_address() {
    int local{42};
    return &local; // после return pointer dangling
}
\end{cppcode}

Это мост к модели lifetime блока~3. \texttt{const} не продлевает жизнь такого
объекта. Sanitizer способен обнаружить некоторые выполненные случаи, но не
обязан диагностировать каждую ошибку lifetime.

\subsection{Pass-by-value: отдельный parameter object}

\begin{cppcode}
void change_copy(int value) {
    value = 99;
}

int original{10};
change_copy(original);
// original всё ещё 10
\end{cppcode}

При pass-by-value параметр --- отдельный локальный объект, инициализированный
из argument. Изменение \texttt{value} не меняет \texttt{original}. Копирование
здесь не случайная цена, а полезная семантика: функция получает собственное
значение.

Для маленьких scalar types вроде \texttt{int}, \texttt{bool} и pointer это
обычно самый ясный выбор. Передача большого \texttt{std::string} или
\texttt{std::vector} по value может копировать элементы и выделять storage.
Точная стоимость зависит от типа и реализации. Иногда копия нужна функции по
смыслу, поэтому правило «reference всегда быстрее» неверно.

\subsection{Reference: alias существующего объекта}

\begin{cppcode}
int score{10};
int& alias{score};
alias = 11; // меняется score
\end{cppcode}

\texttt{T\&} --- reference to \texttt{T}; она --- alias уже
существующего объекта. Reference должна быть привязана при инициализации и
обычно не моделирует optional/null состояние. После привязки assignment через
неё изменяет объект, а не «перенаправляет» reference.

\begin{cppcode}
void add_point(int& score) {
    ++score;
}

int result{10};
add_point(result); // result становится 11
\end{cppcode}

Mutable reference parameter выражает контракт «функция может изменить этот
argument». В месте call нет знака \texttt{\&}, поэтому значимую мутацию важно
подчеркнуть именем: \texttt{add\_point}, \texttt{read\_arguments}.

Reference связана с моделью адреса из блока~3, и реализация часто передаёт её
подобно pointer. Но стандарт задаёт aliasing semantics, а не требует
«скрытый pointer» определённого размера или layout. Pointer уместен, когда
нужно выразить \texttt{nullptr} или менять выбранный адрес; reference --- когда
объект обязателен и alias соответствует контракту. Обе конструкции требуют,
чтобы referent оставался жив.

\subsection{Const reference: read-only interface}

\begin{cppcode}
std::size_t total_characters(
    const std::vector<std::string>& words) {
    std::size_t result{};
    for (const std::string& word : words) {
        result += word.size();
    }
    return result;
}
\end{cppcode}

\texttt{const T\&} связывается с существующим объектом без копирования всего
\texttt{T} и запрещает функции изменять его \emph{через эту reference}.
Поэтому это естественный read-only interface для больших
\texttt{std::string}/\texttt{std::vector}. Компилятор отвергнет изменение
элемента или \texttt{push\_back} через \texttt{words}.

Обещание локально: исходный объект не становится вечным или глобально
immutable; другой non-const alias может его менять. Reference не владеет
объектом и не решает lifetime. Const reference может связываться с temporary и
в некоторых строго определённых контекстах продлевать его lifetime, но это не
универсальный механизм ownership; value categories появятся позже.

\subsection{Четыре роли const в этом блоке}

\begin{cppcode}
const int limit{15};                   // local нельзя изменить
void inspect(int value);               // call interface по value
void inspect_const(const int value);   // тот же тип для caller
void read(const std::string& text);     // read-only alias
const int* view{&limit};                // read-only pointed-to int
\end{cppcode}

Top-level \texttt{const} у value parameter ограничивает только local copy в
definition и не создаёт другую overload. Поэтому его часто опускают в
declaration; в definition можно добавить для защиты реализации. У
\texttt{const T\&} const относится к referent и существенен для call
interface. У \texttt{const int*} нельзя менять pointed-to value через pointer,
но сам pointer можно перенаправить.

\texttt{const} выражает намерение и даёт compiler возможность запретить часть
мутаций. Он не гарантирует thread safety, владение, lifetime или вечность.

\subsection{Return value или output reference?}

\begin{cppcode}
int square(int value) {
    return value * value;
}

void square_out(int value, int& output) {
    output = value * value;
}
\end{cppcode}

Первая форма читается как expression, и результат нельзя забыть
инициализировать внутри функции. Для одного вычисленного результата обычно
предпочтителен return value. Mutable output reference разумна, если главная
цель операции действительно изменить существующий объект. Большое число
output references скрывает направление данных, требует заранее создать
variables и усложняет проверку всех изменений. Более сложные способы возврата
нескольких значений пока не нужны.

\subsection{Overloading и default arguments}

\begin{cppcode}
int twice(int value)       { return value * 2; }
double twice(double value) { return value * 2.0; }

const int a = twice(4);
const double b = twice(2.5);
\end{cppcode}

Overload set содержит функции с одним именем и разными parameter lists.
Compiler выбирает viable candidate и сравнивает требуемые conversions. Если
единственного лучшего кандидата нет, call ambiguous и программа ill-formed.

\begin{cppcode}
void show(int);
void show(double);
// show(1L); // ambiguous: две conversions одного ранга
\end{cppcode}

Return type сам по себе не различает overload: compiler должен выбрать
функцию до использования результата. Перегрузки полезны для одной операции
над близкими типами; одинаковое имя для несвязанных действий скрывает смысл.

Default argument подставляется в месте call:

\begin{cppcode}
void print_line(const std::string& text, int repetitions = 1);
print_line("ready"); // как print_line("ready", 1)
\end{cppcode}

После первого параметра с default все последующие в этом parameter list тоже
должны иметь default. Его обычно указывают в declaration один раз. Это не
новая функция и не runtime-решение; defaults не должны скрывать важные inputs.

\subsection{Первый function pointer}

В блоке~3 pointer мог хранить адрес объекта. Адрес функции тоже можно хранить
и затем вызвать:

\begin{cppcode}
int add(int left, int right) { return left + right; }
int subtract(int left, int right) { return left - right; }

int (*operation)(int, int) = add;
const int first = operation(9, 4); // 13
operation = subtract;
const int second = operation(9, 4); // 5
\end{cppcode}

Тип читается из имени наружу: \texttt{operation} --- pointer на функцию с
двумя \texttt{int} parameters, возвращающую \texttt{int}. Скобки вокруг
\texttt{*operation} существенны. Это простейший callback-like interface:

\begin{cppcode}
int apply(int left, int right, int (*operation)(int, int)) {
    return operation(left, right);
}
\end{cppcode}

Pointer может быть null, поэтому общий interface должен определить это
контрактом или проверить. Синтаксис быстро становится трудным; более удобные
abstractions будут позже вместе с lambdas и \texttt{std::function}. Сейчас
важна идея: поведение тоже можно выбрать через типизированный адрес.

\subsection{Scope, lifetime и повторные вызовы}

Parameter names и locals доступны только в своих scopes. Их automatic
storage duration обычно означает: lifetime начинается при выполнении
declaration во время конкретного call и заканчивается при выходе из блока.
Следующий call создаёт новые parameter/local objects.

\begin{cppcode}
int local_scope_case() {
    const int local{17};
    return local;
}

const int copy = local_scope_case(); // copy жив, local уже нет
// std::cout << local;                // имя вне scope: error
\end{cppcode}

Scope --- область текста, где доступно имя; lifetime --- интервал выполнения,
когда существует объект. Они связаны, но не тождественны. Reference может
покинуть scope исходного имени, но не продлевает lifetime автоматически.
Recursion использует те же правила: каждый call получает свои locals. Её
алгоритмические применения оставим отдельным задачам.

\subsection{Модель стоимости без магии}

Обычный function call может требовать передачи arguments, сохранения точки
возврата и переходов. Оптимизатор может встроить тело (inlining), удалить call
или преобразовать код иначе, если observable behavior сохраняется.

\begin{implementationbox}[Inlining]
Машинное решение зависит от compiler, flags, visibility и target. Keyword
\texttt{inline} не обещает вставку машинного кода; его языковая роль связана
прежде всего с правилами нескольких definitions. И маленькая функция без
keyword может быть встроена, а \texttt{inline}-функция --- нет.
\end{implementationbox}

Копирование большого vector/string по value может копировать элементы и
делать allocation. Const reference часто избегает этой работы, но добавляет
aliasing и зависимость от lifetime, а для маленького \texttt{int} обычно лишь
усложняет interface. Оптимизация способна убрать некоторые копии; утверждение
о времени требует конкретных flags, данных и измерения. Поэтому лаборатория
проверяет семантику, а BENCHMARKS для блока N/A.

\subsection{Эксперимент 1: рефакторинг course\_inventory}

Исходная утилита блока~3 фиксирует поведение: принимает
\texttt{--root}/\texttt{--block}, инспектирует те же 15 пар, печатает ту же
таблицу и возвращает те же классы exit codes. Новая версия лежит в
\texttt{examples/04\_functions\_references\_and\_const/
01\_course\_inventory\_refactored} и не меняет блок~3.

Декомпозиция проведена по задачам: \texttt{read\_arguments} формирует входное
состояние, \texttt{count\_lines} и \texttt{count\_example\_files} измеряют
один ресурс, \texttt{inspect\_block} возвращает одну строку,
\texttt{collect\_rows} строит vector, \texttt{print\_report} агрегирует и
печатает. \texttt{main} теперь выражает верхнеуровневый сценарий.

\texttt{argc}, indexes и selected передаются по value как маленькие значения.
Path, string, массив имён и vector строк отчёта читаются по const reference.
Mutable references \texttt{root}/\texttt{selected} оправданы в parser: его
задача --- заполнить существующее состояние. \texttt{BlockInfo} и vector
естественно возвращаются как values.

\begin{shellcode}
cd examples/04_functions_references_and_const/01_course_inventory_refactored
cmake -S . -B build
cmake --build build
./build/course_inventory_refactored --root ../../.. --block 3
ctest --test-dir build --output-on-failure
\end{shellcode}

Функциональная эквивалентность проверяется одинаковыми arguments и diff после
нормализации только process-specific addresses. PASS --- одинаковый
нормализованный stdout, одинаковые ключевые exit codes и два CTest PASS.

\subsection{Эксперимент 2: function\_contract\_lab}

Каталог \texttt{02\_function\_contract\_lab} содержит один safe executable.
Он проверяет неизменность argument после pass-by-value, изменение через
mutable reference, чтение vector через const reference, return value против
output reference, local lifetime, overloads и выбор add/subtract через
function pointer.

\begin{shellcode}
cd examples/04_functions_references_and_const/02_function_contract_lab
cmake -S . -B build
cmake --build build
./build/function_contract_lab
ctest --test-dir build --output-on-failure
\end{shellcode}

CTest отдельно запускает compiler для трёх ill-formed programs. В
\texttt{const\_mutation.cpp} операция \texttt{+=} требует изменить const
\texttt{std::string}; compiler запрещает это через read-only reference.
Другие cases доказывают, что return type не образует overload и что call может
быть ambiguous. Test считается PASS именно при ненулевом compiler exit code;
реальная diagnostic сохраняется в build.

Unsafe \texttt{dangling\_return} opt-in собирается с ASan/UBSan. PASS требует
ненулевого exit и фактической sanitizer diagnostic на этой toolchain, а не
заранее придуманного текста. Safe executable должен закончить строкой
\texttt{FUNCTION\_CONTRACT\_LAB=PASS} и code 0.

\subsection{Сильные стороны и цена}

Строго типизированные функции дают прозрачную декомпозицию и compile-time
interface. Value/reference/const-reference позволяют явно выбирать копию,
мутацию существующего объекта или read-only alias. Programmer контролирует
форму и потенциальную стоимость передачи данных.

Цена контроля: неверный выбор вызывает лишнюю копию либо неожиданную
мутацию; dangling references/pointers остаются возможны; function pointer
имеет тяжёлый синтаксис; const не решает ownership/lifetime. Compiler проверит
типы, но не определит, хорошо ли программа разделена архитектурно.

\subsection{Проверка понимания}

Нужно уметь:
1) объяснить, когда именованная операция улучшает большой \texttt{main}; 2) различить declaration, definition и call, parameter и argument; 3) описать передачу управления и lifetime locals без требования stack frame; 4) выбрать value, mutable reference или const reference по контракту; 5) объяснить, почему const не означает ownership или вечный объект; 6) предпочесть return value простому output parameter; 7) объяснить overload resolution и два negative cases; 8) прочитать базовый тип function pointer; 9) объяснить dangling return и границы sanitizer; 10) отделить возможный call/copy cost от голословного «быстрее».

\subsection{Источники для сверки}

Function declarations/definitions, calls, return, references,
cv-qualification, overloading и function pointers сверены с C++17 working
draft N4659. Порядок объяснения и исходная задача reuse сопоставлены с MIT
6.096 Lecture~3; references --- с Lecture~5 и релевантными частями Advanced
Topics~II. Их утверждение «reference is just a pointer internally» сознательно
не принято как гарантия стандарта. Современный выбор value/const reference и
C++17 style сопоставлены с \textit{Modern C++ Tutorial}; сборка использует
target-based CMake.
